% titel: Blatt 9, A12: Fällung von Metallsulfiden (S2- über H2S)
% skript: Mortimer Kap. 15 (Löslichkeitsprodukt, Sulfide)

% Teil I: Werkzeuge, kein Lösungsweg.

Ob ein Sulfid ausfällt, zeigt der Vergleich des Ionenprodukts mit dem
Löslichkeitsprodukt.

\begin{werkzeugkasten}[Fällung von Sulfiden]
Zwei Schritte:
\begin{enumerate}
  \item [S$^{2-}$] aus einer gesättigten H$_2$S-Lösung bei gegebenem pH
        berechnen. Für die zweiwertige Säure H$_2$S gilt mit den beiden
        $K_a$-Werten:
        \[
          [S^{2-}] = \frac{K_{a1}K_{a2}\,[\mathrm{H_2S}]}{[\mathrm{H^+}]^2}.
        \]
        ($[H^+]$ aus dem pH, $\mathrm{H_2S}$ gesättigt $=0.1\ M$ hier.)
  \item Ionenprodukt $Q = [M^{2+}][S^{2-}]$ je Metallion bilden und mit dem
        $L$ des Sulfids vergleichen. $Q>L$ $\Rightarrow$ Fällung.
\end{enumerate}
\textit{Skript:} Mortimer Kap. 15 (Löslichkeitsprodukt, Sulfidfällung).
\end{werkzeugkasten}

\begin{achtung}[Darauf achten]
\begin{itemize}
  \item pK$_a$ $\rightarrow$ $K_a = 10^{-pK_a}$.
  \item $[H^+] = 10^{-pH}$ einsetzen, nicht den pH direkt.
  \item Das Ionenprodukt je Metall (Ni$^{2+}$, Co$^{2+}$, Cd$^{2+}$) einzeln
        berechnen und mit dem jeweiligen $L$ vergleichen.
  \item Fällt nur aus, wenn $Q \geq L$.
\end{itemize}
\end{achtung}
